\section{Resumo e Orientações Finais}

% Slide de Matriz de Funcoes
\begin{frame}{Matriz de Decisão: As Quatro Formas de Funções}
    \begin{table}
        \centering
        \scriptsize
        \begin{tabular}{lccl}
            \toprule
            \textbf{Classificação} & \textbf{Entrada} & \textbf{Saída} & \textbf{Exemplo Típico} \\
            \midrule
            Sem Parâm. e Sem Retorno & Não & Não & Banners e menus no console \\
            Com Parâm. e Sem Retorno & Sim & Não & Mensagens de log e alertas \\
            Sem Parâm. e Com Retorno & Não & Sim & Leituras com validação, tokens \\
            Com Parâm. e Com Retorno & Sim & Sim & Cálculos, filtros e buscas \\
            \bottomrule
        \end{tabular}
        \caption{Matriz dos Quatro Tipos de Funções em Programação.}
    \end{table}

    \vspace{0.2cm}
    \begin{block}{Regra Fundamental}
        \begin{itemize}
            \item \textbf{Entrada:} Recebida pelos parâmetros entre os parênteses \texttt{(...)}.
            \item \textbf{Saída:} Devolvida ao código chamador por meio da instrução \texttt{return}.
        \end{itemize}
    \end{block}
\end{frame}

% Slide de Roteiro de Resolucao de Problemas
\begin{frame}{Checklist para Resolução de Algoritmos}
    \begin{columns}
        \begin{column}{0.5\textwidth}
            \begin{block}{1. Análise Prévia}
                \begin{itemize}
                    \item Identifique as entradas necessárias.
                    \item Defina as saídas esperadas.
                    \item Escolha os tipos (\texttt{int}, \texttt{float}, \texttt{list}).
                \end{itemize}
            \end{block}
        \end{column}
        \begin{column}{0.5\textwidth}
            \begin{block}{2. Construção e Teste}
                \begin{itemize}
                    \item Isole a lógica em blocos ou funções.
                    \item Realize o teste de mesa no papel.
                    \item Teste valores de borda ($0$, negativos).
                \end{itemize}
            \end{block}
        \end{column}
    \end{columns}

    \vspace{0.3cm}
    \begin{center}
        \textit{"Primeiro resolva o problema. Depois escreva o código."}
    \end{center}
\end{frame}

% Slide de Fechamento
\begin{frame}{Encerramento e Bons Estudos!}
    \begin{center}
        \LARGE \textbf{Excelente trabalho na revisão!}\\
        \vspace{0.3cm}
        \normalsize Pratique a implementação de cada um dos 25 quesitos.\\
        \vspace{0.5cm}
        \textbf{Prof. Reginaldo Fernandes}\\
        \footnotesize \texttt{reginaldo.fernandes@ifce.edu.br}\\
        \vspace{0.3cm}
        \textbf{IFCE Campus Tauá}\\
        \small Curso Técnico em Redes de Computadores
    \end{center}
\end{frame}
